\documentclass{situationist-preprint}

\title{Cosmic Indifference}
\msid{EXI-06-a}
\msfamily{EXI-06}
\msdiscipline{Meaning, Death and the Absurd}
\mstype{Principal manuscript}
\mskeywords{cosmic indifference; hostility; care; universe; Russell; Hume; Mill}

\begin{document}
\maketitle

\begin{abstract}
Does the universe care about us? The question is often answered in one of two ways: either the cosmos is ordered toward our good, or it is a hostile place that crushes our hopes, as in Russell's picture of a world of blind forces indifferent to human ideals but still somehow arrayed against them. We argue that the universe does not care about us, that it is indifferent, and that indifference is not hostility. Care and hostility are attitudes, and both require a subject whose responses are sensitive to our interests, positively or negatively. The universe, understood as the totality of physical things and the laws they obey, has no such responses. Its outcomes are not tuned to our interests in either direction: the harms it brings are not aimed at us, and the goods are not given to us. The rhetoric of a hostile universe therefore misdescribes the very indifference it is meant to express. The conclusion concerns the universe; it is neutral on whether a being distinct from it cares for us, a question taken up elsewhere.
\end{abstract}

\section{The problem}

\citet{Russell1903} described man as the product of causes that had no prevision of the end they were achieving, and urged a free man's worship of ideals in defiance of an alien and inhuman world; the essay passes, without remark, from the claim that nature is blind to the image of a world that is pitiless and tyrannical. \citet{Mill1874}, in \enquote{Nature}, argued that nature commits, as a matter of daily routine, acts for which men would be imprisoned, and concluded that it could not be a model for conduct. Hume's Philo, closing his catalogue of the miseries of life, judged that the course of nature suggests a blind principle without discernment or parental care \citep[XI]{Hume1779}. And Camus's Meursault, on the night before his execution, opens himself to what he calls the tender indifference of the world \citep{Camus1942}. These authors agree that the universe does not care for us; they differ on whether its want of care is hostility, neutrality or even a kind of consolation. We argue that it is neutrality.

\section{Framework}

\begin{definition}[Responsiveness]\label{d:resp}
A thing is \emph{responsive to our interests} if what it does depends, in a counterfactually robust way, on how its doing so affects our interests: it would do otherwise if our interests were otherwise affected.
\end{definition}

\begin{definition}[Care, hostility, indifference]\label{d:attitudes}
A thing \emph{cares} about us if it is responsive to our interests and its responses tend to promote them; it is \emph{hostile} to us if it is responsive to our interests and its responses tend to set them back; it is \emph{indifferent} to us if it is not responsive to our interests.
\end{definition}

The three conditions are exclusive: care and hostility both require responsiveness and indifference denies it. They are not defined by the outcomes a thing produces. A storm that destroys a village sets back the villagers' interests without being hostile, and a rain that saves a harvest promotes them without caring. What distinguishes care and hostility from mere benefit and harm is that the effect on our interests figures among the conditions on which the thing's behavior depends.

\begin{assumption}[Physical closure]\label{a:phys}
The behavior of the universe, understood as the totality of physical things, is fixed by its prior states and its laws, and the laws make no reference to the interests of any creature.
\end{assumption}

Assumption~\ref{a:phys} is a premise about the content of physical law, not about determinism; the argument goes through if the laws are probabilistic, provided that the probabilities they assign are not themselves sensitive to our interests.

\section{The argument}

\begin{theorem}\label{t:indiff}
The universe is indifferent to us, and hence neither cares about us nor is hostile to us.
\end{theorem}

\begin{proof}
By Assumption~\ref{a:phys}, what the universe does depends only on its prior states and on laws that make no reference to our interests. Fix the prior states: then whether an outcome would set back or promote our interests makes no difference to whether it occurs. The universe is therefore not responsive to our interests in the sense of Definition~\ref{d:resp}, and by Definition~\ref{d:attitudes} it is indifferent. Since care and hostility both require responsiveness, it neither cares about us nor is hostile to us.
\end{proof}

Within the universe there are things that care about us and things that are hostile to us: other people, and perhaps other animals. Their attitudes are realized physically, and the theorem does not deny them. It concerns the universe as a whole and its laws, which are not an agent and are not composed so as to aggregate the attitudes of the agents within them into an attitude of their own.

\begin{corollary}\label{c:rhetoric}
The description of the universe as hostile, pitiless or tyrannical, offered as an expression of its want of care, is false.
\end{corollary}

\begin{proof}
Such a description attributes hostility. By Theorem~\ref{t:indiff}, the universe is not hostile. The description cannot be a more vivid expression of indifference, since by Definition~\ref{d:attitudes} hostility and indifference exclude each other.
\end{proof}

\section{Objections and replies}

\begin{objection}[Harm outweighs benefit]\label{o:harm}
The universe will destroy everything we value, through death, catastrophe and eventually the end of conditions for life. A thing that reliably destroys what we care about is hostile in every sense that matters.
\end{objection}

\begin{reply}
The objection replaces hostility with harmfulness. A harmful thing can be indifferent, as the storm is. Moreover, the universe that will destroy what we value also produced it: the same laws that bring death brought life. If the balance of outcomes made the universe hostile, the existence of the goods would make it caring, and the outcomes cannot settle both questions in opposite ways. Neither does, because neither is a question about outcomes.
\end{reply}

\begin{objection}[Anthropomorphism cuts both ways]\label{o:anthro}
If the universe is not an agent, it is not indifferent either; indifference is an attitude, and to call the universe indifferent is as anthropomorphic as calling it hostile.
\end{objection}

\begin{reply}
Indifference in the sense of Definition~\ref{d:attitudes} is not an attitude but the absence of the responsiveness that attitudes of care and hostility require. A stone is indifferent to us in this sense, without having any attitude. The word misleads when it suggests a turning-away, and that suggestion should be dropped; the claim that remains is the one the theorem establishes.
\end{reply}

\begin{objection}[Theism]\label{o:theism}
If there is a God who created and sustains the universe, as REL-01 argues, the universe is the instrument of a being who cares about us, and the theorem is false.
\end{objection}

\begin{reply}
The theorem concerns the universe, not its ground. If a being distinct from the universe cares about us and acts through or upon the universe, that being cares about us; it does not follow that the universe does, any more than a letter cares about its recipient. Assumption~\ref{a:phys} concerns physical law, and REL-10 argues that the apparent fine-tuning of the physical constants does not support theism. Whether God cares about us, and how that care is compatible with the harms the universe brings, is the subject of REL-02. Our conclusion is neutral on both questions.
\end{reply}

\section{Conclusion}

The universe does not care about us, but it is not against us either. Care and hostility require responsiveness to our interests, and the physical universe, governed by laws that make no reference to them, has none. The language of a hostile universe dramatizes an indifference that is in fact neutral. This bears on the response to the absurd (EXI-04): a world that is not an adversary is not an object of defiance.

\begin{thebibliography}{9}\small
\bibitem[Camus(1942)]{Camus1942} Camus, A. (1942). \emph{L'\'Etranger}. Gallimard.
\bibitem[Hume(1779)]{Hume1779} Hume, D. (1779). \emph{Dialogues Concerning Natural Religion}. Cited by part.
\bibitem[Mill(1874)]{Mill1874} Mill, J. S. (1874). Nature. In \emph{Three Essays on Religion}. Longmans, Green, Reader, and Dyer.
\bibitem[Russell(1903)]{Russell1903} Russell, B. (1903). A free man's worship. \emph{The Independent Review}, 1, 415--424.
\end{thebibliography}

\end{document}
